\begin{exsol}{24}{Tadelis Exercise 4.4}{76--77}{Advanced}{B4}
\asked Under eBay's proxy bidding (the winner pays one increment above the second-highest proxy
bid) with independent private values: (a) argue that bidding above one's value is weakly
dominated by bidding one's value; (b) the same for bidding below; (c) make sense of eBay's
advice to bid one's true value, and say whether you would follow it.

\method Weak dominance (Lecture~1, slide~23; Tadelis Remark, p.~63), by comparing outcomes case
by case in the highest rival proxy bid. As the increment ($\$0.01$) goes to zero, proxy bidding
\emph{is} the second-price auction of item~23. We first argue in that limit, which is the
exercise's intent, and then check what the increment changes.

\solution
Let $v$ be your value, $b$ your proxy bid, $h$ the highest rival proxy bid. In the limit model
you win iff $b > h$ (ties split at random) and pay $h$, whatever $b$ is.

\textbf{(a) Overbidding, $b > v$.} Compare with bidding $v$:
\begin{itemize}
  \item $h < v$: both bids win at price $h$; same payoff $v - h$.
  \item $h > b$: both lose; payoff $0$ for both.
  \item $v \le h \le b$: bid $v$ earns $0$ (it loses, or ties at a price equal to $v$), while
  bid $b$ wins (or ties) at price $h \ge v$, earning $v - h \le 0$. For $v < h < b$ this is
  strictly negative.
\end{itemize}
Bidding $v$ is never worse and is strictly better when $v < h < b$: $b$ is weakly dominated by
$v$. Overbidding only matters when it makes you win at a price above your value.

\textbf{(b) Underbidding, $b < v$.}
\begin{itemize}
  \item $h < b$: both win at price $h$; same payoff.
  \item $h > v$: both lose; same payoff $0$.
  \item $b \le h < v$: bid $v$ wins and earns $v - h > 0$; bid $b$ loses (or ties) and earns
  $0$ (or $(v-h)/2$). Strictly worse.
  \item $h = v$: both earn $0$.
\end{itemize}
So $b$ is weakly dominated by $v$. Underbidding only matters when it makes you lose an auction
you would have won profitably; it never lowers the price you pay.

\textbf{What the increment changes.} With a positive increment the price is
$\min\{b,\ h + \text{increment}\}$ (the winner never pays above his own proxy).
\texttt{check\_C.py} enumerates this rule with bids in cents, a one-cent increment, two bidders,
random tie-breaking and $v = \$10.00$.
\begin{itemize}
  \item (a) holds exactly: $v$ weakly dominates every $b > v$.
  \item (b) holds for every $b \le v - 3$ cents. Bidding $v - 2$ cents is payoff-equivalent to
  bidding $v$. Bidding $v - 1$ cent \emph{weakly dominates} bidding $v$: it ties with probability
  $\tfrac12$ when $h = v - 1$ cent (expected gain half a cent), where bidding $v$ wins at a price of
  exactly $v$ and earns nothing.
\end{itemize}
So the exact statement in (b) holds only up to the size of one increment. Any gain from shading is
below one cent, and it disappears in the limit model the exercise intends.

\textbf{(c) eBay's recommendation.} By (a) and (b), bidding your true value is (weakly) optimal
whatever the other bidders do: you do not need to guess their values, their strategies, or how
many of them there are. Proxy bidding turns eBay's ascending auction into a second-price
auction, and the advice ``enter your true value, no more, no less'' is exactly the weakly dominant
strategy of item~23. I would follow it, with caveats that come from features the model leaves out:
\begin{itemize}
  \item \emph{Dynamics.} Real eBay auctions have a fixed end time and bids arrive over time. Late
  bidding (``sniping'') is common because it gives rivals no time to react, for example to
  incremental bidders who raise their bids step by step. If everyone bids their value once, the
  timing does not matter.
  \item \emph{Values not independent.} For goods with resale or common value, rivals' bids carry
  information about the value (the winner's curse), and ``your value'' must take that into account.
  \item \emph{Increments and shill bidding.} Shading by a few cents is harmless (above); a seller
  bidding on her own item to raise $h$ does not change the advice either, because your bid still
  never affects your price.
\end{itemize}

\interp The general lesson: when your bid affects only \emph{whether} you win and not \emph{what}
you pay, truthful bidding is weakly dominant. The ``weakly'' matters: against many rival profiles
truth-telling only ties with other bids, so it is a safe prescription, not a unique prediction
of what everyone else will do.

\verify \texttt{check\_C.py} checks the limit model as the second-price auction of item~23
(truthful bidding weakly dominates every other grid bid, with one or two rivals) and the
one-cent-increment rule as described above. Tadelis gives no answer in the book.

\answer{(a) Bidding $b > v$ changes the outcome only when $v < h < b$, where it wins at a loss,
so $v$ weakly dominates $b$. (b) Bidding $b < v$ changes the outcome only when $b \le h < v$,
where it gives up a profitable win, so $v$ weakly dominates $b$. (c) Proxy bidding is a
second-price auction, so the true value is a weakly dominant bid, which justifies eBay's advice.
With a positive increment, (b) holds only up to one increment: in cents, $v - 1$ cent weakly
dominates $v$ by half a cent.}
\end{exsol}
